% ======================================================================
%  chapitre3.tex — SPARQL Querying and Analysis
%  \input from main.tex — do not compile standalone.
% ======================================================================

\chapter{SPARQL Querying and Analysis}
\label{chap:queries}

This chapter demonstrates how the EAKG, once deployed in a SPARQL-enabled triplestore such as GraphDB, can be queried and analyzed. We begin with simple exploratory queries, then present a gap analysis technique that leverages the knowledge graph to automate parts of the TOGAF ADM process, and finally describe a set of EA smell detection queries.

% ==================================================================
\section{Simple and Exploratory Queries}
\label{sec:simple-queries}

% ------------------------------------------------------------------
\subsection{Dependency Chain Queries}
\label{subsec:dependency-chains}

Dependency chain queries discover dependencies between elements. These dependencies can be direct or follow a transitive chain using the SPARQL property path operator (\texttt{+}). For example, the following query finds all elements that a given component (here, ``Document Management Back-Up Server'') depends on, traversing both \texttt{archimate:DependencyRelationship} and \texttt{archimate:StructuralRelationship} edges:

\begin{lstlisting}[caption={Dependency chain: elements that a given component depends on}]
PREFIX rdf:  <http://www.w3.org/1999/02/22-rdf-syntax-ns#>
PREFIX rdfs: <http://www.w3.org/2000/01/rdf-schema#>
PREFIX archimate: <http://purl.org/eakg/archimate#>
PREFIX smells: <http://purl.org/eakg/smells#>

SELECT DISTINCT ?component ?service ?serviceName
WHERE {
    ?component rdfs:label "Document Management Back-Up Server"@en .
    ?service (archimate:DependencyRelationship
             |archimate:StructuralRelationship)+ ?component .
    ?service rdf:type archimate:Element ;
             rdfs:label ?serviceName .
}
\end{lstlisting}

All subsequent queries in this chapter and in the Annex assume the same prefixes unless otherwise noted.

Dependency chains can also be traversed in the opposite direction. For example, the query below lists what depends on a given component (here, ``PaaS Services''):

\textit{(full query: see Annex, Listing A.1)}

\begin{table}[H]
\centering
\caption{Dependency chain results for ``Document Management Back-Up Server''}
\label{tab:dependency-results}
\small
\begin{tabularx}{\linewidth}{l l X}
\toprule
\textbf{component} & \textbf{service} & \textbf{service name} \\
\midrule
ex:id-47094 & ex:id-47052 & Home \& Away LAN \\
ex:id-47094 & ex:id-47506 & Homeowner's \& Travel Back Office \\
ex:id-47094 & ex:id-47045 & Data Acquisition Gateway \\
ex:id-47094 & ex:id-47088 & FO General-Purpose Server \\
ex:id-47094 & ex:id-47090 & FO Web Hosting Server \\
ex:id-47094 & ex:id-47508 & Home \& Away Headquarters \\
ex:id-47094 & ex:id-47100 & PaaS Services \\
ex:id-47094 & ex:id-47393 & Project: Data Warehousing and BI \\
ex:id-47094 & ex:id-47501 & Front Office \\
\bottomrule
\end{tabularx}
\end{table}

% ------------------------------------------------------------------
\subsection{Orphan Element Detection}
\label{subsec:orphan-elements}

A useful model quality check is to find elements that have no relationship with any other element---orphan nodes that may have been forgotten during modeling:

\begin{lstlisting}[caption={Detect orphan elements (no incoming or outgoing relationships)}]
SELECT ?element ?elementName
WHERE {
    ?element rdfs:label ?elementName ;
             a archimate:Element .
    FILTER NOT EXISTS {
        ?element archimate:Relationship ?target .
        ?target a archimate:Element .
    }
    FILTER NOT EXISTS {
        ?source archimate:Relationship ?element .
        ?source a archimate:Element .
    }
}
\end{lstlisting}

The query returns 15 orphan elements in the ArchiSurance model, including Product Information, Insurance, Insurance Policy, Commission, Bank, Business Intelligence Team, and several customer-channel elements.

% ------------------------------------------------------------------
\subsection{Most Connected Elements}
\label{subsec:most-connected}

Sorting elements by their total number of relationships (incoming and outgoing) helps identify the most central components of the architecture:

\begin{lstlisting}[caption={Rank elements by total number of relationships}]
SELECT ?element ?elementName
       ((COALESCE(?in,0) + COALESCE(?out,0)) AS ?relations)
WHERE {
    ?element rdfs:label ?elementName .
    ?element a archimate:Element .
    OPTIONAL {
        SELECT ?element (COUNT(DISTINCT ?s) AS ?in)
        WHERE {
            ?s ?rel ?element .
            ?rel rdfs:subPropertyOf* archimate:Relationship .
        }
        GROUP BY ?element
    }
    OPTIONAL {
        SELECT ?element (COUNT(DISTINCT ?t) AS ?out)
        WHERE {
            ?element ?rel ?t .
            ?rel rdfs:subPropertyOf* archimate:Relationship .
        }
        GROUP BY ?element
    }
}
ORDER BY DESC(?relations)
\end{lstlisting}

\begin{table}[H]
\centering
\caption{Eight most connected elements}
\label{tab:most-connected-results}
\small
\begin{tabularx}{\linewidth}{l X r}
\toprule
\textbf{element} & \textbf{element name} & \textbf{relations} \\
\midrule
ex:id-46339 & Target: CRM, Back Office and Data Warehouse Operational & 36 \\
ex:id-46304 & ArchiSurance Back Office Suite & 23 \\
ex:id-47396 & Baseline & 23 \\
ex:id-46637 & Finance & 18 \\
ex:id-46860 & Document Management System & 16 \\
ex:id-46330 & Manage Policies and Claims & 15 \\
ex:id-46918 & Home \& Away Policy Administration & 15 \\
ex:id-46369 & Serve Customers & 14 \\
\bottomrule
\end{tabularx}
\end{table}

% ------------------------------------------------------------------
\subsection{Layer Listings}
\label{subsec:layer-listings}

The following query lists all elements together with their ArchiMate layer classification, exploiting the ontology's class hierarchy:

\textit{(full query: see Annex, Listing A.2)}

Elements can also be listed by the relationships between layers. The cross-layer count below reports the four Application-to-Business type pairs returned by the query.

\textit{(full query: see Annex, Listing A.3)}

\begin{table}[H]
\centering
\caption{Relationships from Application Layer to Business Layer}
\label{tab:cross-layer-results}
\small
\begin{tabularx}{\linewidth}{X X r}
\toprule
\textbf{source type} & \textbf{target type} & \textbf{count} \\
\midrule
archimate:ApplicationComponent & archimate:BusinessFunction & 22 \\
archimate:ApplicationService & archimate:BusinessProcess & 16 \\
archimate:DataObject & archimate:BusinessObject & 3 \\
archimate:ApplicationProcess & archimate:BusinessProcess & 1 \\
\bottomrule
\end{tabularx}
\end{table}

% ==================================================================
\section{Gap Analysis}
\label{sec:gap-analysis}

\subsection{The TOGAF Gap Analysis Technique}
\label{subsec:togaf-gap}

According to the \textit{TOGAF}\textsuperscript{\textregistered} Standard---ADM Technique, Gap Analysis is a technique used to validate an architecture that is being developed \parencite{opengroup2022togaf}. The basic premise is to highlight a shortfall between the Baseline Architecture and the Target Architecture; that is, items that have been deliberately omitted, accidentally left out, or not yet defined.

In other words, Gap Analysis is the process of comparing the current architecture baseline with the new architecture we want to achieve in the future (target) to identify what needs to be added, removed, or modified. A key step in validating the architecture is to identify what has been forgotten.

\noindent\textbf{Suggested Method by the Standard}\par

The standard suggests the following matrix-based method:

\begin{itemize}
    \item Draw a matrix where the Architecture Building Blocks (ABBs) from the baseline architecture are listed per rows, and ABBs present in the target architecture are listed per columns.
    \item Add a last row labeled ``New'' and a last column labeled ``Eliminated.''
    \item ABBs available in both architectures are marked with ``Included'' at the intersecting cell.
    \item ABBs that appear only in the baseline architecture should be reviewed to determine if they were correctly removed or forgotten. If correctly eliminated, mark them in the appropriate ``Eliminated'' cell; otherwise they must be readdressed in the next iteration.
    \item ABBs that appear only in the target architecture are marked in the intersecting cell with the ``New'' row as a gap.
\end{itemize}

When the matrix is completed, ``Eliminated'' and ``New'' marks represent gaps that should be explained---either as correctly eliminated items or as items that need to be addressed by reinstating or developing/procuring the new building block.

% ------------------------------------------------------------------
\subsection{ArchiMate Representation}
\label{subsec:archimate-gap}

In ArchiMate, the ``Implementation and Migration Layer'' includes elements that model the evolution of enterprise architecture through programs and projects \parencite{opengroup2023archimate}. The relevant element types are:

\begin{enumerate}
    \item \textbf{Work Package}: a series of actions identified and designated to achieve specific results within specified time and resource constraints. It has a start and end, and produces goals, outcomes, or deliverables.
    \item \textbf{Deliverable}: represents a result of a Work Package---reports, papers, services, software, physical products, etc.
    \item \textbf{Plateau}: represents a relatively stable state of the architecture that exists during a limited period of time. Plateaus show the enterprise at incremental states reflecting periods of transition between the Baseline and Target Architectures, and allow grouping work packages into portfolios and programs.
    \item \textbf{Gap}: represents a statement of difference between two plateaus. It is associated with two plateaus and represents the difference between them.
\end{enumerate}

% ------------------------------------------------------------------
\subsection{Knowledge Graph Role in Gap Analysis}
\label{subsec:kg-gap}

By transforming the ArchiMate model into a knowledge graph, we can automate the creation of the matrix in the Gap Analysis process. By querying the Knowledge Graph we can extract the list of ``Kept,'' ``Eliminated,'' and ``New'' ABBs. After evaluation by the architect, those elements are marked by linking them to the ``Gap'' node with the appropriate property, to keep track of taken decisions in the knowledge graph.

\noindent\textbf{Queries to List Gap Elements}\par

Three queries extract the gap classification directly from the knowledge graph:

\begin{enumerate}
    \item \textbf{ABBs present in both architectures (Kept).} This query returns the elements aggregated by both the baseline and target plateaus. \textit{(full query: see Annex, Listing A.4)}
    \item \textbf{ABBs present only in the baseline (candidates for elimination).} This query identifies elements that belong to the baseline plateau but are absent from the target plateau. \textit{(full query: see Annex, Listing A.5)}
    \item \textbf{ABBs present only in the target (new gaps).} This query identifies elements introduced in the target plateau that do not exist in the baseline plateau. \textit{(full query: see Annex, Listing A.6)}
\end{enumerate}

\noindent\textbf{Marking Elements After Review}\par

We use project-specific RDF properties to link the ``Gap'' node with each element that presents a gap, which allows us to mark elements with the appropriate decision after review:

\begin{itemize}
    \item \texttt{ex:gapNew} for new ABBs to be produced/procured.
    \item \texttt{ex:gapInclude} for ABBs to be kept, which has a sub-property \texttt{ex:gapImproved} for kept ABBs that require improvement.
    \item \texttt{ex:gapEliminated} for ABBs correctly removed from the architecture.
\end{itemize}

% ------------------------------------------------------------------
\subsection{Example: ArchiSurance Case Study}
\label{subsec:gap-example}

Using the ArchiSurance case study \parencite{jonkers2016archisurance}, we can model the evolution with a ``Gap'' between two plateaus:

\begin{itemize}
    \item The \textbf{baseline plateau} aggregates: Web Portal, Call Center Application, Legal Expense CRM System, General CRM System, Home \& Away Financial Application, Home \& Away Policy Administration, Auto Insurance Application, Legal Expense Back-Office System, Document Management System.
    \item The \textbf{target plateau} aggregates: Web Portal, Call Center Application, General CRM System, Document Management System, ArchiSurance Back-Office Suite, Data Warehousing Solution, MyArchiSurance Mobile App.
\end{itemize}

In the knowledge graph, each plateau is represented with a node that has an ``aggregate'' property linking it to all its included elements, and a ``Gap'' node is associated with both plateau nodes. Using the previous queries, we can list the gap classification without the need to draw the matrix:

\begin{itemize}
    \item \textbf{Kept}: Web Portal, Call Center Application, General CRM System, Document Management System.
    \item \textbf{Eliminated}: Legal Expense CRM System, Home \& Away Financial Application, Home \& Away Policy Administration, Auto Insurance Application, Legal Expense Back-Office System.
    \item \textbf{New}: ArchiSurance Back-Office Suite, Data Warehousing Solution, MyArchiSurance Mobile App.
\end{itemize}

The following combined query classifies elements from the gap analysis according to the TOGAF gap analysis matrix, using concrete identifiers from the ArchiSurance model:

\begin{lstlisting}[caption={Classify elements as Kept, Eliminated, or New in one query}]
# Newly introduced prefix for this query.
PREFIX ex: <http://www.example.org/eakg#>

SELECT ?ABB ?label ?state
WHERE {
    BIND(ex:id-46339 AS ?target)
    BIND(ex:id-47396 AS ?baseline)
    {
        ?baseline archimate:Aggregation ?ABB .
        ?target archimate:Aggregation ?ABB .
        ?ABB rdfs:label ?label .
        BIND("Kept" AS ?state)
    }
    UNION
    {
        ?baseline archimate:Aggregation ?ABB .
        FILTER NOT EXISTS {
            ?target archimate:Aggregation ?ABB .
        }
        ?ABB rdfs:label ?label .
        BIND("Eliminated" AS ?state)
    }
    UNION
    {
        ?target archimate:Aggregation ?ABB .
        FILTER NOT EXISTS {
            ?baseline archimate:Aggregation ?ABB .
        }
        ?ABB rdfs:label ?label .
        BIND("New" AS ?state)
    }
}
\end{lstlisting}

After reviewing each element, we mark it with the appropriate property by inserting triples with the ``Gap'' node. The following SPARQL Update query shows how decisions are recorded:

\textit{(full query: see Annex, Listing A.7)}

% ------------------------------------------------------------------
\subsection{Further Exploration: Ontology-Based EA Evolution}
\label{subsec:gap-further}

A small ontology to model enterprise evolution is proposed in \parencite{silva2017evolution}, taking advantage of computer agents to aid communication and analysis of EA model evolution. The key concepts of this ontology are:

\begin{itemize}
    \item An \textbf{IT Project} acts as an enabler of EA evolution by applying a set of \textbf{Transformations}.
    \item A Transformation can be decomposed into smaller \textbf{Changes}, each altering the life-cycle of an EA Description (element or relationship).
    \item To keep track of motivation and interested parties, the ontology introduces \textbf{Stakeholder} (who has one or more drivers) and \textbf{Driver} (an internal or external motivation/condition that makes the change necessary).
    \item This model keeps information about how the EA was changed, when those changes happened, who involved those changes, and why.
\end{itemize}

The authors proposed a set of Description Logic (DL) queries to answer questions that help decision-making and facilitate change analysis. In our case, DL queries can be replaced with equivalent SPARQL queries into the knowledge graph.

\begin{samepage}
For example, to find what transformations are applied by an IT project named ``Back-Office System Integration,'' the proposed DL query:

\begin{lstlisting}[language=, basicstyle=\small\ttfamily, frame=single]
Transformation and
  applied_by some
    (IT_Project and
     hasName value "Back-Office_System_Integration")
\end{lstlisting}
\end{samepage}

can be expressed as the following SPARQL equivalent:

\textit{(full query: see Annex, Listing A.8)}

% ==================================================================
\section{EA Smells Detection Queries}
\label{sec:smells}

Enterprise Architecture Smells are recurring patterns in an EA model that indicate potential design problems \parencite{salentin2020smells}. The following subsections present SPARQL queries that detect and tag several categories of smells, leveraging the RDFS reasoning capabilities of the triplestore. For each smell, a detection \texttt{SELECT} identifies the affected resources and a tagging \texttt{INSERT} materializes a \texttt{smells:hasSmell} triple on each one. These queries were tested on the demonstration model from \parencite{salentin2020smells}, which was designed to exhibit each smell.

% ------------------------------------------------------------------
\subsection{Structural and Connectivity Smells}
\label{subsec:structural-smells}

\noindent\textbf{Cyclic Dependency}\par

\textbf{The Concept:} Elements should not be stuck in an endless loop of depending on one another. This behavior is illustrated in the demonstration model of \parencite{salentin2020smells}.

\textbf{The Detection Logic:} We target the \texttt{archimate:DependencyRelationship} hierarchy and scan the graph globally to match unbroken chains of these dependency edges. If traversing that chain eventually leads right back to the exact same starting node, a cycle is detected.

\textit{(full query: see Annex, Listing A.9)}

Each element in the cycle is tagged with \texttt{smells:CyclicDependency}. \textit{(tagging query: see Annex, Listing A.14)}

\noindent\textbf{Shared Persistency}\par

\textbf{The Concept:} Higher-level business components should not bypass the application layer to talk directly to the same backend database. This behavior is illustrated in the demonstration model of \parencite{salentin2020smells}, where it is treated together with Strict Layers Violation rather than as a separately defined smell.

\textbf{The Detection Logic:} We isolate \texttt{archimate:PassiveStructureElement} nodes, which represent data. We then look for instances where two or more distinct elements from the \texttt{archimate:BusinessLayer} have an \texttt{archimate:Access} edge pointing directly to that exact same data node.

\textit{(full query: see Annex, Listing A.10)}

The shared data node is tagged with \texttt{smells:SharedPersistency}. \textit{(tagging query: see Annex, Listing A.15)}

\noindent\textbf{Strict Layers Violation}\par

\textbf{The Concept:} The architecture should follow a strict top-down hierarchy (Business $\rightarrow$ Application $\rightarrow$ Technology). Skipping an intermediate layer by establishing a direct link between the Business layer and the Technology layer breaks this core design principle. As a best practice, no direct architectural relationships of any kind or direction should bridge these layers directly without traversing the application layer.

\textbf{The Detection Logic:} We search for any direct relationship where one connected node belongs to the \texttt{archimate:BusinessLayer} and the other belongs to the \texttt{archimate:TechnologyLayer}. By checking both directions, we confirm that the \texttt{archimate:ApplicationLayer} has been entirely bypassed, regardless of how the relationship was originally oriented in the model.

\begin{lstlisting}[caption={Detect strict layers violation}]
SELECT DISTINCT ?nodeA ?nodeB ?relation ?direction
WHERE {
    {
        ?nodeA a archimate:BusinessLayer .
        ?nodeB a archimate:TechnologyLayer .
        ?nodeA ?relation ?nodeB .
        BIND("Business_to_Tech" AS ?direction)
    }
    UNION
    {
        ?nodeA a archimate:TechnologyLayer .
        ?nodeB a archimate:BusinessLayer .
        ?nodeA ?relation ?nodeB .
        BIND("Tech_to_Business" AS ?direction)
    }
    ?relation rdfs:subPropertyOf* archimate:Relationship .
    FILTER(STRSTARTS(STR(?relation),
           "http://www.example.org/eakg#id"))
}
\end{lstlisting}

The offending relationship edge is tagged with \texttt{smells:StrictLayersViolation}. \textit{(tagging query: see Annex, Listing A.16)}

% ------------------------------------------------------------------
\subsection{Service and Interaction Smells}
\label{subsec:service-smells}

\noindent\textbf{Chatty Service}\par

\textbf{The Concept:} A service should not be overly communicative, firing off a massive number of granular requests to a vast array of other services. \parencite{salentin2020smells} state that a service related to a high number of other services may indicate this smell.

\textbf{The Detection Logic:} We target service nodes (\texttt{archimate:ApplicationService}, \texttt{archimate:BusinessService}) and count the total number of unique elements they connect to. If that count exceeds a defined acceptable threshold, the service is flagged.

\begin{lstlisting}[caption={Detect chatty services (threshold: 10 connections)}]
SELECT ?service ?label
       (COUNT(DISTINCT ?connectedElement) AS ?connectionCount)
WHERE {
    ?service a ?serviceType .
    FILTER (?serviceType IN (archimate:ApplicationService,
                             archimate:BusinessService))
    ?service rdfs:label ?label .
    {
        ?service ?relation ?connectedElement .
    }
    UNION
    {
        ?connectedElement ?relation ?service .
    }
    ?relation rdfs:subPropertyOf* archimate:Relationship .
    FILTER(STRSTARTS(STR(?relation),
           "http://www.example.org/eakg#id"))
}
GROUP BY ?service ?label
HAVING (COUNT(DISTINCT ?connectedElement) >= 10)
\end{lstlisting}

The overly communicative service node is tagged with \texttt{smells:ChattyService}. \textit{(tagging query: see Annex, Listing A.17)}

\noindent\textbf{Data Service}\par

\textbf{The Concept:} A service should execute logic. If it only touches data objects and nothing else, it is just a weak data-access point. This smell is described in \parencite{salentin2020smells} as a service that only references data objects without relating to other components.

\textbf{The Detection Logic:} We isolate service nodes and evaluate their \emph{outgoing} relationships (where the service acts as the source). We identify services that have outgoing edges pointing to \texttt{archimate:PassiveStructureElement} nodes, but apply a strict filter to ensure they have \emph{zero} outgoing edges pointing to any \texttt{archimate:ActiveStructureElement} or \texttt{archimate:BehaviorElement}. Incoming edges (other components calling the service) are expected and are ignored by this filter.

\textit{(full query: see Annex, Listing A.11)}

The data-only service node is tagged with \texttt{smells:DataService}. \textit{(tagging query: see Annex, Listing A.18)}

% ------------------------------------------------------------------
\subsection{Model Quality Smells}
\label{subsec:quality-smells}

\noindent\textbf{Duplication}\par

\textbf{The Concept:} The model should not contain multiple distinct elements that share the exact same name. Sharing a name alone is treated as a smell, regardless of the elements' specific types. A similar case is illustrated in the demonstration model of \parencite{salentin2020smells}, based on name similarity rather than exact matching.

\textbf{The Detection Logic:} We compare the \texttt{rdfs:label} (the name) across the graph. If we find two or more distinct nodes (different IRIs) sharing the exact same label, they are flagged as duplicates.

\textit{(full query: see Annex, Listing A.12)}

Each element sharing its label with at least one other element is tagged with \texttt{smells:Duplication}. \textit{(tagging query: see Annex, Listing A.19)}

\noindent\textbf{Unnecessary Documentation}\par

\textbf{The Concept:} Model documentation should be concise. Huge walls of text clutter the repository. This smell is illustrated in the demonstration model of \parencite{salentin2020smells}, but no general threshold is stated.

\textbf{The Detection Logic:} We extract the \texttt{rdfs:comment} string literal for every element in the graph, calculate the character length of that string, and flag comments that exceed a specific character limit (e.g., 400 characters).

\textit{(full query: see Annex, Listing A.13)}

Any model resource exceeding the character limit is tagged with \texttt{smells:UnnecessaryDocumentation}. \textit{(tagging query: see Annex, Listing A.20)}

\subsection{Utility Functions}
\label{subsec:utility-functions}

The following query lists the elements tagged with a selected smell; replace
\texttt{smells:CyclicDependency} with any smell URI as needed:

\begin{lstlisting}[caption={View tagged elements for a specific smell}]
SELECT ?element ?label WHERE {
    ?element smells:hasSmell smells:CyclicDependency .
    OPTIONAL { ?element rdfs:label ?label . }
}
\end{lstlisting}

A corresponding query for resetting a selected smell tag is provided in Annex,
Listing A.21.
